\chapter{Why the whole path has one endpoint value}\label{ch:path}

Local marginals change while an action unfolds. The path identity shows why adding those changes gives exactly the finite EBU already defined. It also separates improvement between endpoints from service, exposure and resource use accumulated through time, giving each a clear place in an EBU account.


\section{Many small changes, one complete comparison}
A pipe can transfer water steadily or in pulses. During the movement the local marginal changes: the source becomes less full and the destination less empty. If we keep using the first marginal, we will count later units as though the earlier ones never arrived. The integral solves precisely this problem. It follows the changing marginal along the entire represented path.

There are two objects here. The finite definition compares states, $E=V(x_0)-V(x_1)$. The path theorem explains when a continuous calculation gives that same number. The theorem does not define a new kind of EBU or add a charge for traveling between the endpoints. It connects the finite account to local differential information.

Think of measuring elevation along a walk. Adding all upward and downward elevation changes gives the difference between the final and initial heights. A long detour changes the sequence of readings, but not that height difference. Distance walked and effort expended are different quantities. In an EBU application, pumping energy and wear likewise need their own physical representation; they cannot be inferred from the endpoint change of an inventory potential.

\section{The assumptions that make the connection exact}
\begin{theorem}{Fixed-field path identity}\label{thm:path}
Let $V$ be a single-valued $C^1$ function on an open neighbourhood of an admissible piecewise $C^1$ path $\gamma:[0,1]\to\mathcal X$, with $\gamma(0)=x_0$ and $\gamma(1)=x_1$. Hold the potential and its parameters fixed throughout the comparison. Then
\begin{equation}\label{eq:path-main}
 E=V(x_0)-V(x_1)
   =-\int_0^1\nabla V(\gamma(s))^{\mathsf T}\gamma'(s)\,ds.
\end{equation}
The same statement holds intrinsically on a smooth constrained state manifold, using its differential $dV$ along a piecewise smooth path.
\end{theorem}
\begin{proof}
On each smooth segment the chain rule gives
$\frac{d}{ds}V(\gamma(s))=\nabla V(\gamma(s))^{\mathsf T}\gamma'(s)$.
Integrating gives the potential at that segment's end minus the potential at its start. Summing over segments cancels the intermediate endpoints. Negating the result gives the stated sign convention.
\end{proof}

The proof is short because the important work is in the declaration. We have supplied a state function, the path lies in its domain, and the derivative is regular enough for the fundamental theorem of calculus. We have not started with an arbitrary vector field and tried to prove that a potential exists. That distinction removes a frequently misplaced topological condition: \emph{no simply connected domain is required when the single-valued potential is already supplied}. A hole in the domain does not invalidate the chain rule for a path on which $V$ is defined.

If instead someone gives only a vector field and asks whether it is a gradient, global compatibility becomes a separate problem. Vanishing local curl alone can fail on a domain with holes. That is not the problem solved by defining EBU through an already declared $V$.

\section{What the local differential says}
Write
\begin{equation}
 dE=-dV=-\nabla V(x)^{\mathsf T}dx.
\end{equation}
This is an infinitesimal relation along the fixed field. It says how a tiny represented displacement contributes to the full comparison. It does not say that a finite displacement can be priced by one derivative sample. For an edge increment $dx=S_e\,dq$, it becomes $dE=f_e(x)\,dq$, where $f_e=-\nabla V^{\mathsf T}S_e$.

A useful analogy is a meter whose reading changes while a tank fills. Multiplying the first rate by the full duration works only when the relevant rate stays constant. Integrating changing readings is not an optional correction. It is the calculation needed to describe the whole event.

The path parameter $s$ is not necessarily clock time. Replacing it by a smooth increasing parameter changes both the velocity and the integration element, leaving their product and the integral unchanged. Thus the mathematical account need not depend on how quickly an interpolation is traversed. A physical consequence that does depend on speed, such as heating or spoilage, requires a state and action model that actually represents it.

\section{The same endpoints do not mean the same operation}
Two pumps can transfer the same water while using different electricity. A potential over tank stocks assigns equal values to their matching water endpoints. A potential over the fuller carrier state can distinguish their energy consequences. The identity remains exact in either case; the boundary determines which comparison it makes.

Time introduces another distinction. Along a physical trajectory $x(t)$ on $[0,T_f]$,
\begin{equation}\label{eq:path-versus-exposure}
 -\int_0^{T_f}\nabla V(x(t))^{\mathsf T}\dot x(t)\,dt
 =V(x(0))-V(x(T_f)),\qquad
 \mathcal A=\int_0^{T_f}V(x(t))\,dt.
\end{equation}
The first expression integrates the \emph{rate of change} of the potential. The second integrates its \emph{level}. They are different functions of a trajectory and have different units: $\mathcal A$ includes an extra unit of time. Two deliveries reaching the same endpoint can have equal EBU while one leaves the clinic short of water for longer.

Feedback makes this distinction especially valuable. Overshoot, settling time, peak shortage and cumulative exposure describe different aspects of service. A controller's oscillation can affect all of them without producing an extra transaction or changing the endpoint definition. A delivery policy can therefore be compared on both restoration and continuity of service, with neither measure hidden inside the other.

For EBU, path independence gives a concrete protection: changing the bookkeeping route between fixed complete endpoints cannot manufacture additional total value. The separate time description gives a complementary benefit: a system can value prompt, reliable service without mislabelling duration as another endpoint EBU. These two capabilities support clearer comparisons of maintenance and route choices.

\section{An integral checked against the endpoints}
For the familiar two-cell transfer, $V(q)=16-4q+q^2/4$ and $f(q)=4-q/2$. A four-unit movement therefore has
\begin{equation}
 E=\int_0^4(4-q/2)\,dq=16-4=12.
\end{equation}
Direct evaluation gives the same answer from $V(0)=16$ and $V(4)=4$. A two-stage physical execution gives seven for the first two units and five for the next two. The changing contributions sum to twelve because their baselines join.

The agreement is more informative than a numerical coincidence. If a software quote returns sixteen, it probably kept the first field value for the whole movement. If it returns fourteen by adding two same-baseline two-unit quotes, it probably treated two counterfactual beginnings as one physical history. The endpoint identity gives an independent way to recognize both mistakes.

\section{A path that is declared is not necessarily observed}
A straight segment between the endpoints is often convenient. On a convex stock domain it stays within that domain. But it is a physical trajectory only if the declared mechanism supplies that timing and simultaneous flow. Endpoint measurements alone cannot prove that every intermediate point occurred.

For the total, every admissible path satisfying the theorem gives the same result. For a decomposition among children or actors, the chosen path can matter. We will therefore calculate the joint outcome and its interaction structure before treating a path split as an attribution rule. A convenient interpolation must remain visible as a convention.

\section{Regularity is part of the theorem}
The $C^1$ assumptions above are sufficient and easy to check. A weaker formulation can use absolute continuity of $V\circ\gamma$ together with its almost-everywhere chain rule and an integrable derivative. Merely having a derivative almost everywhere is not enough: singular functions can change while their derivative is zero almost everywhere. Appendix~\ref{app:calculus} explains the corresponding issue for repeated finite differences.

A jump retains its endpoint EBU whenever both potentials exist, even when its actual mechanism has no smooth-path description. The finite account therefore survives the boundary of the differential representation. The next chapter evaluates it exactly for a quadratic, including the complete curvature correction.
